\section{One filling increment in every degree}\label{sec:ordinary}

We combine fixed-degree shortening with the filtered cut--join identity.
The cut removes a degree-dependent amount of length, while its join with
a sphere power restores that length up to one fixed difference.  The
result is a filling increment independent of both degree and level.

\begin{theorem}[Uniform ordinary death]\label{thm:ordinary-death}
For every smooth Riemannian metric on $S^n$, $n\ge3$, there is a constant
$D>0$ such that, for every integer $j\ge0$ and every real $a>0$,
\begin{align}
 \ker\bigl(H_j(X_a)\longrightarrow H_j(\Lambda)\bigr)
 &\subseteq
 \ker\bigl(H_j(X_a)\longrightarrow H_j(X_{a+D})\bigr),
 \label{eq:ordinary-absolute}\\
 \ker\bigl(H_j(X_a,M)\longrightarrow H_j(\Lambda,M)\bigr)
 &\subseteq
 \ker\bigl(H_j(X_a,M)\longrightarrow H_j(X_{a+D},M)\bigr).
 \label{eq:ordinary-relative}
\end{align}
The same constant works at regular and nonregular levels.  No assumption
on the number of prime geodesic images is required.
\end{theorem}

\begin{proof}
We first prove the absolute assertion for the length filtration
$\mathcal A_a$, then pass to square-root energy and to the relative pair.

\paragraph{Choosing the degree reduction.}
Let $d=n-1$ and let $\alpha,C_H$ be the constants of
Proposition~\ref{prop:sphere-powers}.  For sufficiently large $j$ put
\begin{equation}\label{eq:ordinary-degrees}
 m=\left\lfloor\frac{j-2}{2d}\right\rfloor,
 \qquad q=2md,\qquad p=q-d,\qquad N=n+q.
\end{equation}
Then
\begin{equation}\label{eq:ordinary-remainder}
 2\le r:=j-q\le2d+1=2n-1,
 \qquad j>q+1=n+p,
 \qquad N>n+p.
\end{equation}
Choose an integer $J>2n$ so large that, whenever $j\ge J$, these choices
have $m\ge1$, $p>n+1$, and $L>0$ in
\eqref{eq:cj-thresholds}.  Take $S,T,l,L$ as in that equation.  Write
\begin{equation}\label{eq:ordinary-gap}
 C_{\mathrm{gap}}=l-L
   =\frac{2n-1+2C_H}{\alpha}+4,
 \qquad
 D_* =\max_{0\le s<J}D_s,
\end{equation}
where $D_s$ is furnished by Proposition~\ref{prop:fixed-degree-death}.
All the remainder degrees in \eqref{eq:ordinary-remainder}, and the
degree $n$, occur in this finite maximum.

\paragraph{Filling a high-degree null cycle.}
Let $x$ be a finite globally null $j$-cycle in $\mathcal A_a$, with
$j\ge J$.  Its cut $D_Tx$ is a globally null $r$-cycle in
$\mathcal A_{a-L+1}$ by Lemma~\ref{lem:cj-operations}.
If $a-L+1>0$, Proposition~\ref{prop:fixed-degree-death} gives a chain
$W$ such that
\[
 \partial W=D_Tx,
 \qquad\operatorname{supp}W\subset\mathcal A_{a-L+1+D_r}.
\]
If $a-L+1\le0$, the cut has no surviving term and we take $W=0$.

Fix a fundamental constant-loop cycle $c$.  By
Proposition~\ref{prop:absolute-unit}, the cycle $D_TS-c$ is globally
null in degree $n$ and is supported in
$\mathcal A_{C_{\mathrm{gap}}+1}$.  Fixed-degree filling gives a chain
$V$ with
\begin{equation}\label{eq:ordinary-correction}
 \partial V=D_TS-c,
 \qquad
 \operatorname{supp}V\subset
 \mathcal A_{C_{\mathrm{gap}}+1+D_n}.
\end{equation}
This length bound is independent of $j$, although the chain $V$ itself
may depend on the chosen power.

Choose the join radius after these finite length bounds have been
specified.  Lemma~\ref{lem:cj-operations} and
Theorem~\ref{thm:filtered-interchange} permit one choice for the
interchange homology and both joined fillings.  Refinement comparisons
are included in the same bounds.  The relevant chains have the following
supports:
\begin{align}
 \mu(W,S)&\text{ lies in }
   \mathcal A_{a+C_{\mathrm{gap}}+D_*+2},
   \label{eq:ordinary-joined-filling}\\
 \mu(V,x)&\text{ lies in }
   \mathcal A_{a+C_{\mathrm{gap}}+D_*+2}.
   \label{eq:ordinary-joined-correction}
\end{align}
The first assertion is omitted when $W=0$; otherwise it follows by
adding the input bounds $a-L+1+D_r$ and $l$ and the join allowance $1$.
The second follows from \eqref{eq:ordinary-correction} in the same way.
Both chains have degree $j+1$.

The degree inequalities \eqref{eq:ordinary-remainder} are precisely
the two hypotheses needed for the interchange homology.  It gives
\[
 \mu(D_Tx,S)\simeq\mu(D_TS,x)
 \quad\text{in }\mathcal A_{a+C_{\mathrm{gap}}+4}.
\]
The chain $\mu(W,S)$ has boundary $\mu(D_Tx,S)$, while
$\mu(V,x)$ has boundary
$\mu(D_TS,x)-\mu(c,x)$.  Finally
\eqref{eq:cj-chain-unit} joins $\mu(c,x)$ to $x$ in
$\mathcal A_a$.  Adding these chains and the interchange homology
therefore fills $x$.

Thus one permissible length increment is
\begin{equation}\label{eq:ordinary-D}
 D=C_{\mathrm{gap}}+D_*+6.
\end{equation}
It exceeds every bound just used.  For $j<J$, the conclusion follows
directly from Proposition~\ref{prop:fixed-degree-death}, since
$D_j\le D_*<D$.  We have proved that every globally null finite cycle
in $\mathcal A_a$, in any degree, fills in $\mathcal A_{a+D}$.
The constants in \eqref{eq:ordinary-D} depend only on the metric and
$n$, and every construction used above is independent of a finite-prime
hypothesis.

\paragraph{Returning to square-root energy.}
A cycle $x$ in $X_a$ is in $\mathcal A_a$ because $\ell\le F$.
Fill it by the length assertion and apply $\Psi_1$ to that filling.
Lemma~\ref{lem:arclength} gives $F(\Psi_1\gamma)=\ell(\gamma)$, so
the new filling lies in $X_{a+D}$ and has boundary $\Psi_1x$.
The normalization prism from $x$ to $\Psi_1x$ lies entirely in $X_a$
by \eqref{eq:shortening-energy}.  Adding this prism proves
\eqref{eq:ordinary-absolute} with the same $D$.  Strict sublevel
inequalities are preserved throughout; no approximation by regular
levels is involved.

\paragraph{Passing to constants-relative homology.}
Let $i:M\to\Lambda$ include constants and let
$\ev_0:\Lambda\to M$ be evaluation at zero.  Since
$\ev_0\circ i=\operatorname{id}_M$, the chain map
\[
 P=\operatorname{id}-i_\#(\ev_0)_\#
\]
annihilates $C_*(M)$ and preserves $C_*(X_a)$ for every $a>0$.
If $z$ represents a relative cycle in $(X_a,M)$, then
$\partial z\in C_*(M)$, so $Pz$ is an absolute cycle representing
the same relative class.  If $z$ is globally null relative to $M$,
there are finite chains $U$ in $\Lambda$ and $c'$ in $M$ with
$\partial U=z+c'$.  Applying $P$ gives $\partial(PU)=Pz$.
Hence $Pz$ is globally null absolutely.  By
\eqref{eq:ordinary-absolute} it fills in $X_{a+D}$, and passage to
the relative chain quotient proves \eqref{eq:ordinary-relative}.
This argument includes degree zero and uses exactly the same increment.
\end{proof}

The restriction $a>0$ ensures that the constant loops belong to the
relative pair.  For nonpositive $a$ the open absolute sublevel is empty.
Theorem~\ref{thm:ordinary-death} supplies the degree-uniform control of
ordinary transient classes needed for the reflection-equivariant argument
below.
